\subsection{Output Units}
The output units are the last layer of the network, and their choice depends on the type of problem we are trying to solve.
\begin{itemize}
    \item \textbf{Linear Units}: used for regression problems, where the output is simply a linear combination of the last hidden layer. 
    For a given input $\vct{x}$, the output of the linear layer can be expressed as:
    \begin{equation*}
        y = \vct{\theta}^T \vct{x} + \theta_0
    \end{equation*}
    Linear units usually do not saturate, which makes them effective for gradient-based optimization methods, 
    as they do not suffer from the vanishing gradient problem.
    
    \item \textbf{Softmax Units}: used for multi-class classification problems. 
    A softmax layer takes the unnormalized output of the last hidden layer and converts it into a probability distribution over the classes.
    For a given input $\vct{x}$, the output of the softmax layer can be expressed as:
    \begin{equation*}
        \hat{y}_i = \frac{\exp(z_i)}{\sum_{j=1}^{K} \exp(z_j)}
    \end{equation*}
    where $z_i$ is the output of the last hidden layer for class $i$, and $K$ is the total number of classes.

    The $exp$ function ensures that the output is always positive, and the normalization term ensures that the outputs sum to 1, making them naturally interpretable as probabilities.
\end{itemize}
